\documentclass[11pt]{article}
\usepackage{fontspec}
\setmainfont{Latin Modern Roman}
\usepackage{amsmath,amssymb,mathtools,unicode-math}
\setmathfont{Latin Modern Math}
\usepackage{xurl,hyperref}
\hypersetup{colorlinks=true,urlcolor=blue,pdftitle={Problem statement}}
\newcommand{\sref}[2]{\hyperlink{source:#1:#2}{[S#2]}}
\newcommand{\cataloguescope}[1]{\paragraph{Mathematical question.}#1}
\begin{document}
\section{Research Problems in Function Theory — Problem 2.85}\label{statement}

\subsection{Definitions and mathematical statement}

\paragraph{Shared notation.}$\mathbb N=\{1,2,\ldots\}$, and $\mathbb Z,\mathbb Q,\mathbb R,\mathbb C$ denote the integers, rationals, reals and complex numbers. Any other conventions are specified locally.


Let $P_\lambda(z)=\lambda(z+z^2)$, where $|\lambda|=1$ and $\lambda$ is not a root of unity. There is a unique formal series $h_\lambda(z)=z+O(z^2)$ satisfying $P_\lambda(h_\lambda(z))=h_\lambda(\lambda z)$. Let $R(\lambda)\in[0,\infty]$ be its radius of convergence and put $m=\sup_\lambda R(\lambda)$. A periodic cycle of least period $q$ is repelling if $|(P_\lambda^{\circ q})'(z)|>1$ at a point of the cycle.
\paragraph{Statement recorded in the supplied source.}
(a) Determine $m$ exactly, or approximate it with absolute error at most $10^{-10}$. (b) Is the supremum attained for some $\lambda=e^{2\pi i\alpha}$ with $\alpha$ a real quadratic irrational? (c) If $R(\lambda)=0$, does every disk $\{|z|<\varepsilon\}$, $\varepsilon>0$, contain an entire repelling periodic cycle of $P_\lambda$? \sref{statement}{2}
\subsection{Short English statement}
Find the maximal linearization radius for the quadratic family, determine whether a quadratic irrational rotation attains it, and ask whether failure of linearization forces repelling cycles arbitrarily close to the fixed point.
\subsection{Sources}
\begin{description}
\item[\hypertarget{source:statement:1}{S1}] ULAM.AI and the UnsolvedMath Contributors, UnsolvedMath: A Curated Collection of Open Mathematics Problems, catalogue entry. CC BY 4.0. \url{https://huggingface.co/datasets/ulamai/UnsolvedMath}.
\item[\hypertarget{source:statement:2}{S2}] W. K. Hayman and E. F. Lingham, \emph{Research Problems in Function Theory}, arXiv:1809.07200v2 (2018), Problem 2.85. \url{https://arxiv.org/abs/1809.07200}.
\end{description}
\label{end:statement}
\end{document}
